% ============================================================
% APPENDIX E — A Real-World ADLC Framework
% ============================================================
\chapter{A Real-World ADLC Framework: Analysis of a Professional Contract for AI Agents}\index{ADLC!framework}\index{SKILL.md}
\label{app:adlc-framework}

\section*{Foreword}

Throughout the book you have learned how to write \texttt{\_CONTEXT.md} files, define constraints, and structure your interaction with the AI agent. This appendix shows how the same principles are applied in a \textbf{production-grade} framework made of multiple modules and specialized skills.

The framework analyzed here is included in the repository under \texttt{example\_contracts/}. It is not a thought experiment: it represents a practical, structured contract for real software projects.

% ---------------------------------------------------------
\section{Framework Architecture}

The framework is organized across \textbf{three hierarchical levels}, just like the Progressive Disclosure pattern introduced in Chapter~3.

\begin{lstlisting}[language={},title={Framework structure}]
Repository/
├── .github/
│   ├── copilot-instructions.md          ← Universal contract
│   └── copilot_modules/                 ← Specialized modules
│       ├── 00_MODE_EN.md
│       ├── 00_CONTEXT_TEMPLATE_EN.md
│       ├── 01_CORE_RULES_EN.md
│       ├── 02_DISCOVERY_ANALYSIS_EN.md
│       ├── 03_DESIGN_EN.md
│       ├── 04_IMPLEMENTATION_EN.md
│       ├── 05_VERIFICATION_RELEASE_EN.md
│       ├── 06_OPS_EN.md
│       ├── 07_SPECIAL_LANES_EN.md
│       ├── 08_PROMPT_LIBRARY_EN.md
│       ├── 09_CODEBASE_ANALYSIS_EN.md
│       ├── 10_DOCUMENTATION_EN.md
│       ├── 11_BUGFIX_PLAYBOOK_EN.md
│       ├── 12_OPERATING_GUIDE_EN.md
│       ├── 13_ARCH_BACKEND_EXAMPLE_EN.md
│       ├── SECURITY_CONSTRAINTS_LIBRARY_EN.md
│       └── PERFORMANCE_CONSTRAINTS_LIBRARY_EN.md
│
├── .copilot/
│   └── skills/
│       ├── SKILL_ANALYSIS_EN.md
│       ├── SKILL_DESIGN_EN.md
│       ├── SKILL_REACT_EN.md
│       ├── SKILL_FLUTTER_EN.md
│       ├── SKILL_API_DESIGN_EN.md
│       ├── SKILL_DATABASE_EN.md
│       ├── SKILL_SECURITY_EN.md
│       └── SKILL_OPS_EN.md
│
├── ProjectA/
│   ├── .copilot/
│   │   ├── instructions.md
│   │   └── domain.md
│   ├── _CONTEXT.md
│   └── PROGRESS.md
│
└── ProjectB/
    ├── .copilot/
    ├── _CONTEXT.md
    └── PROGRESS.md
\end{lstlisting}

\subsection{The separation principle}

\begin{center}
\begin{tabularx}{\textwidth}{l l l X}
\toprule
\textbf{Scope} & \textbf{Location} & \textbf{Managed by} & \textbf{Purpose} \\
\midrule
Universal framework & \texttt{.github/} & Architecture team & Immutable rules shared across projects \\
Project rules & \texttt{.copilot/} & Project team & Stack- and domain-specific overrides \\
Session state & \texttt{\_CONTEXT.md} & Human + agent & Current phase, active task, constraints, commands \\
\bottomrule
\end{tabularx}
\end{center}

This architecture solves a practical problem: how to keep consistency across many projects without copying the same instructions into every repository.

% ---------------------------------------------------------
\section{The Universal Contract: \texttt{copilot-instructions.md}}

The main file acts as a \textbf{router}. It defines the essential protocols and delegates details to specialized modules. Four elements are especially important.

\subsection{Bootstrap Protocol}
At the start of each new session, the agent:
\begin{enumerate}
  \item discovers the codebase structure,
  \item locates the nearest \texttt{\_CONTEXT.md},
  \item loads project instructions and skills,
  \item produces a structured context block with phase, stack, constraints, and commands.
\end{enumerate}

\subsection{Operating Modes}

\begin{center}
\begin{tabularx}{\textwidth}{l l l X}
\toprule
\textbf{Mode} & \textbf{Active checks} & \textbf{Speed} & \textbf{Typical use} \\
\midrule
STANDARD & SEC/PERF/test/doc checks & Normal & Day-to-day development \\
FAST & Minimal critical checks & 2$\times$ & Prototypes and spikes \\
AUDIT & Full checks plus compliance & Slow & Pre-production review \\
RAPID & Only basic sanity checks & 5$\times$ & Emergency fixes \\
\bottomrule
\end{tabularx}
\end{center}

\subsection{Risk Classification and Mandatory STOP}

The framework formalizes the same logic you saw in the book:
\begin{lstlisting}[language={}]
LOW RISK + SMALL SCOPE:
→ Execute directly and report when done

MEDIUM RISK or MEDIUM SCOPE:
→ Propose the change and request confirmation

HIGH RISK or LARGE SCOPE:
→ Present a detailed plan; explicit approval required
\end{lstlisting}

Mandatory STOP situations include database schema changes, API contract changes, deletion of data or code, authentication changes, production deploys, and secrets management.

\subsection{Confidence Tagging}
Every substantial technical answer ends with a confidence tag and an evidence basis:
\begin{lstlisting}[language={}]
---
AI CONFIDENCE: FACT
Basis: Function signature verified in src/auth/handler.ts#L42
\end{lstlisting}

\begin{notebox}
This is the professional extension of the Confidence Tagging idea: not just \emph{FACT} or \emph{ASSUMPTION}, but also the exact evidence used to justify the claim.
\end{notebox}

% ---------------------------------------------------------
\section{SDLC Modules: One Agent for Each Phase}

A key innovation of the framework is modularization \textbf{by development phase}. Instead of loading every rule in every session, the system activates only the relevant modules.

\begin{center}
\begin{small}
\begin{tabularx}{\textwidth}{X l X}
\toprule
\textbf{Input pattern} & \textbf{Detected phase} & \textbf{Module loaded} \\
\midrule
``I need...'' / ``new feature'' & 0--1: Discovery & \texttt{02\_DISCOVERY\_ANALYSIS\_EN.md} \\
``architecture'' / ``ADR'' / ``tech stack'' & 2: Design & \texttt{03\_DESIGN\_EN.md} \\
``implement'' / ``code'' / ``write'' & 3: Implementation & \texttt{04\_IMPLEMENTATION\_EN.md} \\
``test'' / ``QA'' / ``is it ready?'' & 4--5: Verification & \texttt{05\_VERIFICATION\_RELEASE\_EN.md} \\
``production error'' / ``incident'' & 6: Ops & \texttt{06\_OPS\_EN.md} \\
\bottomrule
\end{tabularx}
\end{small}
\end{center}

\subsection{The Design module}
It contains stack evaluation templates, ADR structures, API contract patterns, and EPIC-to-task decomposition.

\subsection{The Implementation module}
It enforces the \textbf{Pseudocode Rule}: for tasks with complex logic or more than roughly 50 lines of algorithmic code, outline the logic first, then implement.

\subsection{The Ops module}
It covers incident triage, root cause analysis, post-mortems, and operational runbooks.

% ---------------------------------------------------------
\section{Interaction Workflow: The Human-Agent Protocol}

The full working cycle has five phases:
\begin{enumerate}
  \item \textbf{Initial setup}: the human prepares the project structure and base context
  \item \textbf{Session start}: the agent runs the bootstrap protocol
  \item \textbf{Task assignment}: by task file, by epic, or freeform
  \item \textbf{Work cycle}: implementation with checkpoints every 3--5 meaningful actions
  \item \textbf{Session end}: update of \texttt{\_CONTEXT.md} and \texttt{PROGRESS.md}
\end{enumerate}

\subsection{Control commands}

\begin{center}
\begin{tabularx}{\textwidth}{l X}
\toprule
\textbf{Command} & \textbf{Effect} \\
\midrule
\texttt{@stop} & Immediate halt \\
\texttt{@explain} & Explain the reasoning behind the last action \\
\texttt{@undo} & Revert the last change \\
\texttt{@alternatives} & Propose alternative solutions \\
\texttt{@confidence} & Declare FACT / INFERRED / ASSUMPTION \\
\texttt{@context-update} & Update \texttt{\_CONTEXT.md} and \texttt{PROGRESS.md} \\
\texttt{@checkpoint} & Save progress and propose the next step \\
\bottomrule
\end{tabularx}
\end{center}

% ---------------------------------------------------------
\section{Specialized Components}

\subsection{Parallel full-stack workflow}
For features that require synchronized frontend and backend changes, the framework defines explicit sync points:
\begin{lstlisting}[language={}]
SYNC POINT 1: Contract Agreement
→ OpenAPI spec finalized, shared DTOs agreed

SYNC POINT 2: Integration Ready
→ Backend operational, frontend ready for real integration

SYNC POINT 3: E2E Ready
→ Both layers integrated and testable end-to-end
\end{lstlisting}

\subsection{Constraint libraries}
Two dedicated libraries collect reusable patterns:
\begin{itemize}
  \item \texttt{SECURITY\_CONSTRAINTS\_LIBRARY\_EN.md}
  \item \texttt{PERFORMANCE\_CONSTRAINTS\_LIBRARY\_EN.md}
\end{itemize}

\subsection{SKILL files: competencies for every phase}

\begin{center}
\begin{small}
\begin{tabularx}{\textwidth}{l l X}
\toprule
\textbf{SKILL} & \textbf{Phase} & \textbf{Competency} \\
\midrule
\texttt{SKILL\_ANALYSIS\_EN.md} & 0--1 & Requirements, user stories, threat model \\
\texttt{SKILL\_DESIGN\_EN.md} & 2 & Stack evaluation, ADR, task design \\
\texttt{SKILL\_REACT\_EN.md} & 3--4 & React frontend implementation \\
\texttt{SKILL\_FLUTTER\_EN.md} & 3--4 & Flutter mobile implementation \\
\texttt{SKILL\_API\_DESIGN\_EN.md} & 3--4 & REST API design and contracts \\
\texttt{SKILL\_DATABASE\_EN.md} & 3--4 & PostgreSQL, Prisma, migrations \\
\texttt{SKILL\_SECURITY\_EN.md} & 3--4 & OAuth~2.0, JWT, OWASP practices \\
\texttt{SKILL\_OPS\_EN.md} & 6 & Incident handling and runbooks \\
\bottomrule
\end{tabularx}
\end{small}
\end{center}

This is Progressive Disclosure applied to \emph{skills}: only the relevant expertise is loaded when it is needed.

% ---------------------------------------------------------
\section{From the Book to the Framework: A Mapping}

\begin{center}
\begin{small}
\begin{tabularx}{\textwidth}{l c X}
\toprule
\textbf{Book concept} & \textbf{Ch.} & \textbf{Framework implementation} \\
\midrule
\texttt{\_CONTEXT.md} & 3 & Standardized context card template \\
SEC/PERF constraints & 3, 14 & Reusable constraint libraries \\
Confidence Tagging & 3, 14 & Tag plus mandatory evidence basis \\
Risk Classification & 14 & Risk matrix with Mandatory STOP \\
Pseudocode Rule & 10 & Mandatory for complex tasks \\
Context Engineering & 3 & Three-level Progressive Disclosure \\
Multi-agent pattern & 16 & Modules loaded per phase \\
Checkpointing & 10, 16 & Protocol every 3--5 actions \\
Testing and security & 14 & Checklists and release gates \\
Deploy discipline & 15 & Rollback plan and release governance \\
\bottomrule
\end{tabularx}
\end{small}
\end{center}

% ---------------------------------------------------------
\section{Lessons from the Real World}

\subsection{Framework immutability matters}
The universal rules belong in \texttt{.github/}. Project teams customize \texttt{.copilot/}, not the core framework.

\subsection{Context cost must be managed}
Each module declares its approximate size and context cost, helping teams avoid saturation.

\subsection{Priority hierarchy must be explicit}
\begin{enumerate}
  \item \texttt{\_CONTEXT.md} --- highest priority
  \item \texttt{.copilot/instructions.md}
  \item \texttt{.copilot/skills/*.md}
  \item \texttt{.copilot/domain.md}
  \item \texttt{.github/copilot\_modules/} --- lowest priority
\end{enumerate}

\subsection{Control commands act like an operating system}
Commands such as \texttt{@stop}, \texttt{@explain}, and \texttt{@checkpoint} are not decoration; they are the human control surface for the agent.

% ---------------------------------------------------------
\section{Beyond the Textual Contract: Infrastructure Governance}

Textual constraints are powerful, but they still rely on the model respecting them. Enterprise systems increasingly add \textbf{runtime enforcement} through sandboxing and policy-based controls.

A prominent example is NVIDIA's \textbf{NemoClaw}, which combines textual guidance with kernel-level policy enforcement. In practice this means:
\begin{itemize}
  \item network access can be limited to approved endpoints,
  \item file system access can be restricted to authorized directories,
  \item sensitive data can be routed through local models rather than public cloud systems.
\end{itemize}

\begin{center}
\begin{tabularx}{\textwidth}{l X X}
\toprule
& \textbf{Textual Contract (ADLC)} & \textbf{Infrastructure Governance (NemoClaw)} \\
\midrule
Mechanism & Natural-language rules & YAML policies at runtime \\
Enforcement & The agent \emph{should} comply & The agent \emph{cannot} violate \\
Cost & Near zero & Dedicated infrastructure \\
Ideal for & Small teams and individual projects & Enterprise and regulated environments \\
Limitation & Vulnerable to prompt injection & Higher complexity and cost \\
\bottomrule
\end{tabularx}
\end{center}

\begin{notebox}
For the kind of projects covered in this book, the textual ADLC contract is usually more than sufficient. Infrastructure-level governance becomes necessary only when you deal with regulated data, highly sensitive environments, or large fleets of autonomous agents.
\end{notebox}

% ---------------------------------------------------------
\section*{Conclusion}

This framework demonstrates that ADLC is not an academic curiosity. It is a practical operating model for real teams, real repositories, and real delivery constraints. The path from a simple \texttt{\_CONTEXT.md} to a multi-module professional framework is not a leap --- it is a progressive staircase built from the same core ideas you have already learned.
